\section*{अध्याय \ref{ch:g11:reaction-energy} --- अभिक्रियाओं की ऊर्जा: दहन और आबंध ऊर्जाएँ}

\begin{solution}{exo:g11:reaction-energy:1}
लकड़ी का \omterm{def:g4:what-a-fire-needs:burning}{जलना}: \omterm{def:g11:reaction-energy:exothermic}{ऊष्माक्षेपी}। बर्फ़ का पिघलना: \omterm{def:g11:reaction-energy:exothermic}{ऊष्माशोषी} (वह अपने परिवेश से ऊष्मा
लेती है, यद्यपि यह \omterm{def:g7:chemical-reactions:reaction}{रासायनिक अभिक्रिया} नहीं है)। शीत पैक:
\omterm{def:g11:reaction-energy:exothermic}{ऊष्माशोषी}। हाथ गरम करने वाली थैली: \omterm{def:g11:reaction-energy:exothermic}{ऊष्माक्षेपी}। धूप में शर्करा बनाना:
\omterm{def:g11:reaction-energy:exothermic}{ऊष्माशोषी} (ऊर्जा प्रकाश से आती है)।
\end{solution}

\begin{solution}{exo:g11:reaction-energy:2}
$Q = 3.0 \times 890.7 = \qty{2.7e3}{kJ}$, लगभग \qty{2.7}{MJ}।
\end{solution}

\begin{solution}{exo:g11:reaction-energy:3}
\omterm{def:g7:chemical-reactions:reactant}{अभिकारकों} में। अंतर परिवेश को मुक्त होता है, ज़्यादातर
ऊष्मा के रूप में।
\end{solution}

\begin{solution}{exo:g11:reaction-energy:4}
$n = 100 / 46.0 = \qty{2.17}{mol}$; $Q = 2.17 \times 1367.6 =
\qty{2.97e3}{kJ}$, लगभग \qty{3.0}{MJ}।
\end{solution}

\begin{solution}{exo:g11:reaction-energy:5}
ऐसे आबंधों के एक \omterm{def:g10:the-mole:amount}{मोल} को तोड़ने के लिए आवश्यक ऊर्जा, \omterm{def:g7:atoms-and-molecules:molecule}{अणु} और \omterm{def:g7:atoms-and-molecules:atom}{परमाणु}
गैस अवस्था में मानकर। साझा युग्म दोनों \omterm{def:g7:atoms-and-molecules:atom}{परमाणुओं} को साथ बाँधे रखता है: उन्हें
अलग खींचना उस आकर्षण के विरुद्ध काम करना है, जिसमें ऊर्जा लगती है।
\end{solution}

\begin{solution}{exo:g11:reaction-energy:6}
टूटे: $2 \times 436 + 498 = \qty{1370}{kJ}$; बने:
$4 \times 464 = \qty{1856}{kJ}$; $E_r \approx 1370 - 1856 =
\qty{-486}{kJ}$: \omterm{def:g11:reaction-energy:exothermic}{ऊष्माक्षेपी}।
\end{solution}

\begin{solution}{exo:g11:reaction-energy:7}
एथेनॉल \ce{CH3-CH2-O-H} में 5 \ce{C-H}, 1 \ce{C-C}, 1 \ce{C-O} और 1
\ce{O-H} हैं। टूटे: $5 \times 415 + 345 + 350 + 464 + 3 \times 498 =
\qty{4728}{kJ}$। बने: $4 \times 741 + 6 \times 464 = \qty{5748}{kJ}$।
$E_r \approx \qty{-1020}{kJ}$, परिमाण में मापे गए \qty{1367.6}{kJ} से
लगभग एक चौथाई कम।
\end{solution}

\begin{solution}{exo:g11:reaction-energy:8}
प्रोपेन: $2219.2 / 0.0440 = \qty{5.04e4}{kJ/kg} = \qty{50.4}{MJ/kg}$।
मेथेन: $890.7 / 0.0160 = \qty{55.7}{MJ/kg}$। दोनों आलेख के अनुसार।
\end{solution}

\begin{solution}{exo:g11:reaction-energy:9}
$Q = 1500 \times 4.18 \times 85 = \qty{5.33e5}{J} = \qty{533}{kJ}$;
$n = 533 / 890.7 = \qty{0.598}{mol}$, अर्थात $0.598 \times 16.0 =
\qty{9.6}{g}$ मेथेन; आधी ऊर्जा नष्ट होने पर, \qty{19}{g}।
\end{solution}

\begin{solution}{exo:g11:reaction-energy:10}
हाइड्रोजन > मेथेन > प्रोपेन > ऑक्टेन > एथेनॉल।
$48.0 / 142.9 = \qty{0.34}{kg}$ हाइड्रोजन।
\end{solution}

\begin{solution}{exo:g11:reaction-energy:11}
टूटे: $946 + 3 \times 436 = \qty{2254}{kJ}$; बने:
$6 \times 390 = \qty{2340}{kJ}$; $E_r \approx \qty{-86}{kJ}$: थोड़ा
\omterm{def:g11:reaction-energy:exothermic}{ऊष्माक्षेपी}।
\end{solution}

\begin{solution}{exo:g11:reaction-energy:12}
पानी: $250 \times 4.18 \times 20 = \qty{2.09e4}{J} = \qty{20.9}{kJ}$।
एथेनॉल: $1.20 / 46.0 = \qty{0.0261}{mol}$, जो
$0.0261 \times 1367.6 = \qty{35.7}{kJ}$ मुक्त कर सकता था। दक्षता
$20.9 / 35.7 \approx \qty{59}{\percent}$। हानियाँ: \omterm{def:g4:air-a-mixture-of-gases:air}{वायु} द्वारा ले जाई गई ऊष्मा,
डिब्बे और स्टैंड को गरम करने वाली ऊष्मा, \omterm{def:g8:combustion-and-fuels:complete}{अपूर्ण दहन} (कालिख),
कुछ एथेनॉल का बिना जले वाष्पित होना।
\end{solution}

\begin{solution}{exo:g11:reaction-energy:13}
\ce{2C8H18 + 25O2 -> 16CO2 + 18H2O}: प्रति \qty{5470}{kJ} कार्बन डाइऑक्साइड के 8 \omterm{def:g10:the-mole:amount}{मोल},
\qty{352}{g}, इसलिए प्रति मेगाजूल $352 / 5.470 = \qty{64.4}{g}$।
\ce{C3H8 + 5O2 -> 3CO2 + 4H2O}: प्रति \qty{2219.2}{kJ} \qty{132}{g},
इसलिए प्रति मेगाजूल \qty{59.5}{g}। मेथेन,
\qty{49.4}{g} के साथ, सबसे कम मुक्त करती है।
\end{solution}

\begin{solution}{exo:g11:reaction-energy:14}
सारणी औसत \omterm{def:g11:reaction-energy:bond-energy}{आबंध ऊर्जाएँ} देती है, जबकि कार्बन डाइऑक्साइड के \ce{C=O} आबंध
औसत से अधिक प्रबल हैं; मापे गए मान द्रव पानी के लिए हैं,
जिसका संघनन गैस अभिक्रिया से अधिक ऊर्जा मुक्त करता है;
और विधि हर आबंध को उसके पड़ोसियों से स्वतंत्र मानती है। अनुमान फिर भी उपयोगी है:
वह बिना किसी माप के सही चिह्न (\omterm{def:g11:reaction-energy:exothermic}{ऊष्माक्षेपी}) और परिमाण की
सही कोटि देता है।
\end{solution}

\begin{solution}{exo:g11:reaction-energy:15}
टूटे: $4 \times 464 = \qty{1856}{kJ}$; बने: $2 \times 436 + 498 =
\qty{1370}{kJ}$; $E_r \approx \qty{+486}{kJ}$: \omterm{def:g11:reaction-energy:exothermic}{ऊष्माशोषी}, इसलिए ऊर्जा
देनी पड़ती है। बाद में हाइड्रोजन को \omterm{def:g4:what-a-fire-needs:burning}{जलाने} पर यह ऊर्जा वापस मिलती है:
हाइड्रोजन कहीं और उत्पन्न ऊर्जा को संचित करती है, स्वयं कोई ऊर्जा नहीं बनाती।
\end{solution}

\begin{solution}{pb:g11:reaction-energy:1}
\textbf{1.} \ce{CH4 + 2O2 -> CO2 + 2H2O}।

\textbf{2.} \ce{C2H6O + 3O2 -> 2CO2 + 3H2O}।

\textbf{3.} \ce{2C8H18 + 25O2 -> 16CO2 + 18H2O}।

\textbf{4.} \ce{2H2 + O2 -> 2H2O}।

\textbf{5.} हाइड्रोजन।

\textbf{6.} \qty{16.0}{g/mol}, \qty{46.0}{g/mol}, \qty{114.0}{g/mol},
\qty{2.0}{g/mol}।

\textbf{7.} \si{kJ} प्रति \si{kg} में, फिर \si{MJ/kg} में: मेथेन
$890.7 / 0.0160$, \qty{55.7}{MJ/kg}; एथेनॉल $1367.6 / 0.0460$,
\qty{29.7}{MJ/kg}; ऑक्टेन $5470 / 0.1140$, \qty{48.0}{MJ/kg}; हाइड्रोजन
$285.8 / 0.0020$, \qty{142.9}{MJ/kg}।

\textbf{8.} हाइड्रोजन > मेथेन > ऑक्टेन > एथेनॉल।

\textbf{9.} हाइड्रोजन बहुत हल्की गैस है: सामान्य दाब पर उसका एक किलोग्राम
विशाल आयतन घेरता है, इसलिए उसे भारी
उच्च-दाब टंकियों में ठूँसकर भरना पड़ता है।

\textbf{10.} मेथेन: 4 \ce{C-H}; डाइऑक्सीजन: 2 \ce{O=O} टूटते हैं। कार्बन
डाइऑक्साइड: 2 \ce{C=O}; पानी: $2 \times 2 = 4$ \ce{O-H} बनते हैं।

\textbf{11.} $4 \times 415 + 2 \times 498 = \qty{2656}{kJ}$।

\textbf{12.} $2 \times 741 + 4 \times 464 = \qty{3338}{kJ}$।

\textbf{13.} $E_r \approx 2656 - 3338 = \qty{-682}{kJ}$, जबकि मापा गया
\qty{-890.7}{kJ} है: परिमाण में लगभग \qty{23}{\percent} कम।

\textbf{14.} औसत \omterm{def:g11:reaction-energy:bond-energy}{आबंध ऊर्जाएँ} (कार्बन डाइऑक्साइड का \ce{C=O}
औसत से अधिक प्रबल है); माप में द्रव पानी, अनुमान में
गैसें।

\textbf{15.} $1000 / 890.7 = \qty{1.123}{mol}$।

\textbf{16.} $1.123 \times 44.0 = \qty{49.4}{g}$।

\textbf{17.} एथेनॉल: $1000 / 1367.6 = \qty{0.731}{mol}$, जिससे
कार्बन डाइऑक्साइड के \qty{1.462}{mol}, \qty{64.3}{g}, मिलते हैं। ऑक्टेन:
$1000 / 5470 = \qty{0.183}{mol}$, जिससे \qty{1.463}{mol}, \qty{64.4}{g}।

\textbf{18.} मेथेन: उसमें प्रति कार्बन \omterm{def:g7:atoms-and-molecules:atom}{परमाणु} सबसे अधिक हाइड्रोजन \omterm{def:g7:atoms-and-molecules:atom}{परमाणु} हैं (4,
ऑक्टेन के 2.25 के मुक़ाबले), और हाइड्रोजन के \omterm{def:g4:what-a-fire-needs:burning}{जलने} से बिना कार्बन डाइऑक्साइड के
ऊर्जा मिलती है।

\textbf{19.} इस पर कि हाइड्रोजन कैसे बनाई जाती है: नवीकरणीय स्रोतों की
बिजली से पानी से बनी हाइड्रोजन लगभग कोई कार्बन डाइऑक्साइड मुक्त नहीं करती; मेथेन
से बनी हाइड्रोजन कार्बन डाइऑक्साइड को निकास-नली के बजाय कारख़ाने में
मुक्त करती है।

\textbf{20.} मेथेन प्रति मेगाजूल लगभग \qty{49}{g} कार्बन डाइऑक्साइड
मुक्त करती है।
\end{solution}
